% ==========================================================================
% CHAPTER 6 -- CONCLUSION AND FUTURE PLAN
% ==========================================================================

\chapter{Conclusion and Future Plan}
\label{ch:conclusion}

This chapter closes the book. Section~\ref{sec:ch6_conclusion} restates the
problem and the contribution of the work, Section~\ref{sec:limitations}
states honestly what the work does not cover, and Section~\ref{sec:future}
proposes specific next steps.


\section{Conclusion}
\label{sec:ch6_conclusion}

The problem addressed by this thesis was the unreliable and often inflated
reporting of ECG arrhythmia classification performance, combined with the
difficulty of running accurate classifiers on the small microcontrollers
found in wearable monitors. To answer it, a
hybrid CNN--BiLSTM--attention model and a lightweight multi-scale
feature-augmented network (MSFT-Net) were designed, trained, and compared
under a strictly record-disjoint, inter-patient partition of the MIT-BIH
data, following the ANSI/AAMI EC57 class grouping.

Each objective of Section~\ref{sec:objectives} was met. The first
objective, designing the two architectures with rhythm, morphological, and
template-similarity auxiliary features, was achieved in
Chapter~\ref{ch:methodology}: MSFT-Net (167{,}090 parameters) attained
89.45\% accuracy and 87.25\% macro-$F_1$ under the inter-patient protocol,
significantly outperforming the 3.4$\times$ larger recurrent hybrid
(82.00\% accuracy, 73.26\% macro-$F_1$, $p \approx 0$ by McNemar test),
while the same architecture family reached 98.59\% accuracy under the
permissive intra-patient six-category protocol. The second objective, honest
evaluation under both protocols with statistical significance testing, was
met in Chapter~\ref{ch:results}, and it produced the central finding of the
thesis: the evaluation protocol changes the reported performance by roughly
nine to ten points and even reverses the architectural conclusion, since
the recurrent stage helps under the subject-mixed split but actively harms
inter-patient generalisation. The third objective, component analysis and
faithful explainability, was met by the twelve-arm ablation, which showed
that the auxiliary feature vector is the engine of the model (removing it
costs 0.2819 macro-$F_1$) and that the attribution maps, while plausibly
localised to the QRS complex, fail the deletion-curve faithfulness test and
must be treated as qualitative. The fourth objective, INT8 deployment on an
STM32 microcontroller, was met with a measured per-beat inference latency
of 185.057~ms, a 288.16~KB model occupying 56.28\% of flash and 75.97\% of
RAM, and the full inter-patient accuracy preserved on-device.

The contribution of this work is threefold: (i) a protocol-correct,
statistically disciplined evaluation of two competing ECG classification
architectures under identical splits, quantifying for the first time in a
single controlled study how much apparent performance is a by-product of
the assessment regime; (ii) a compact, feature-augmented convolutional
architecture (MSFT-Net) with a staged, evidence-based treatment of the
resistant fusion class; and (iii) a demonstrated deployment path from
deep learning model to real-time INT8 inference on actual
microcontroller hardware, with measured rather than asserted
deployability.


\section{Limitations}
\label{sec:limitations}

The following limitations qualify the conclusions above.

First, the $Q$ class rests on very few patients: its test support is a
single held-out paced record, so its sensitivity and positive predictivity
are de facto single-patient statistics whose real uncertainty is the
between-patient spread, not the nominal binomial confidence intervals.
Second, the CHF category of the intra-patient study comprises ECG segments
derived from a database that
reflects only hospitalised patients with congestive heart failure, recorded
at a different sampling frequency than the arrhythmia database; the
source-database probe confirmed a learnable database fingerprint in the
integrated corpus, so the near-perfect separation of the CHF-derived
category under the
intra-patient protocol must not be interpreted as a cardiac-morphology
finding or as evidence of a clinically generalisable CHF signature. Third, robustness was measured only with synthesised Gaussian
noise and simulated baseline wander; this is not equal to validation on an
independently collected ambulatory recording, so the level of real-world
sensor and motion artefact to which the results generalise is untested.
Fourth, the latency of the smallest MCU model (Table~\ref{tab:mcu}) was
characterised on a general-purpose host, not on target embedded hardware,
so its timing reflects the computational footprint rather than a
demonstration of real-time on-device operation; the full-accuracy MSFT-Net
was, however, profiled directly on the STM32. Fifth, both architectures
were evaluated on data from a few public databases and a single
record-disjoint partition; a completely out-of-hospital patient population
would be a more rigorous test of generalisation than any within-database
split can provide. Finally, the calibration and decision-threshold
experiment showed that thresholds tuned on the validation partition do not
transfer to the test partition, so any operating-point tuning performed for
deployment must be verified on data from a more representative deployment
population than an internal split can provide.


\section{Future Plan}
\label{sec:future}

Several concrete next steps follow from the limitations above:

\begin{enumerate}
  \item Extend the evaluation to an external, independently collected
        ambulatory cohort (for example a wearable single-lead database), so
        that the robustness and generalisation claims are validated on real
        sensor and motion artefacts rather than synthetic corruption.
  \item Increase the $Q$-class support by adding more paced records from
        other public databases, so that the between-patient statistics of
        the $Q$ class become meaningful rather than single-patient.
  \item Improve the faithfulness of the attribution maps by training with
        explanation objectives (for example attention or saliency
        regularisation), and re-test with deletion curves and
        intersection-over-union against cardiologist-annotated intervals.
  \item Investigate the hypothesis, raised by the reversal of the
        recurrent-stage benefit between protocols, that recurrent
        architectures trained on multi-subject data preferentially learn
        patient-identity-correlated temporal structure, for example by
        probing the learned recurrent representations for patient identity.
  \item Add on-device probability calibration with abstention so the
        wearable can withhold low-confidence classifications instead of
        emitting unreliable decisions, and validate the chosen operating
        point on a deployment-representative population.
  \item Extend the deployment work to continuous multi-hour operation on
        the STM32 (real-time acquisition, buffering, and communication
        stacks alongside inference), and measure battery life and power
        consumption in a wearable form factor.
  \item Retrain the pipeline on multi-lead and 12-lead data and on
        additional arrhythmia classes beyond the five AAMI classes, to move
        the system towards broader clinical screening coverage.
\end{enumerate}

These steps are the natural continuation of the work: each one follows
directly from a limitation identified in Section~\ref{sec:limitations} or
from a finding of Chapter~\ref{ch:results}, and together they outline the
path from the present prototype to a clinically validated, wearable cardiac
monitoring system.